% Provenance (never printed): sources/book2/BRIEFS.md, g10-04. Market demand, the law of demand,
% what moves the curve. Real numbers: figdata/grade-10/04-demand/ (BLS average prices and CPI,
% EIA, WDI population); the Pink Sheet fallback of the sync (S2): US gasoline price from BLS.
% Jensen and Miller (2008): ledger g10:jensen-miller-2008. Brenland numbers are invented.
\chapter{Demand}\label{ch:g10:demand}

In 1978 a gallon of regular gasoline cost an American driver \usd{3.22} in today's money. Three years
later it cost \usd{4.75}, almost half as much again, after two shocks to the world supply of oil. Americans
did not stop driving to work, heating their homes or flying. But they bought smaller cars, drove a little
less, insulated their houses, and switched power stations from oil to coal. By 1983, the United States
consumed \qty{23.8}{barrels} of oil products per person a year, against \qty{30.9}{barrels} in 1978: almost a
quarter less. Thirty years later the same thing happened again, more slowly: after the price rises of 2005--2012,
oil use per person fell by a sixth.
% figdata: oil_us (1978 3.22, 30.91; 1981 4.75; 1983 23.78; 2005 25.69; 2012 4.98, 21.46)

When a price rises, people buy less of the good. This is the most robust finding in economics, and it is
the subject of this chapter. But the opening chart also shows how hard it is to read off a single number:
over the same years, incomes grew, cars changed and the country's population aged. A demand curve is a
statement about one cause at a time, and the record mixes them all.

\begin{openingchart}{The price of regular gasoline in the United States, in 2024 dollars per gallon (top), and
the consumption of refined oil products per person, in barrels per year (bottom), 1976--2024. One barrel is
about 159 litres.}{US Bureau of Labor Statistics, average retail prices and CPI-U all items; US Energy
Information Administration, international energy statistics; World Bank, World Development Indicators
(population)}
\begin{tikzpicture}
\begin{groupplot}[group style={group size=1 by 2, vertical sep=0.9cm, x descriptions at=edge bottom},
  omtime, width=0.88\linewidth, height=4.6cm, xmin=1975, xmax=2025, xtick={1980,1985,...,2020}]
\nextgroupplot[ymin=1.5, ymax=5.5, ylabel={dollars per gallon}]
\addplot[omThm, thick, mark=*, mark size=0.8pt] table[x=year, y=price] {figdata/grade-10/04-demand/out/oil_us.csv};
\nextgroupplot[ymin=18, ymax=32, ylabel={barrels per person}]
\addplot[omDef, thick, mark=*, mark size=0.8pt] table[x=year, y=bpp] {figdata/grade-10/04-demand/out/oil_us.csv};
\end{groupplot}
\end{tikzpicture}
\end{openingchart}

\begin{recall}
From the grades 6--9 volume: the \emph{quantity demanded} of a good is the amount buyers want to buy at a given
price; a \emph{demand table} lists it for several prices, and the \emph{demand curve} draws it, with the price on
the vertical axis. A demand curve can be a straight line, $Q^d = a - b\,p$. When something other than the price
changes, the whole curve moves: a \emph{shift} of the curve. From \cref{ch:g10:economic-models}: a prediction
about the price holds \emph{\omterm{def:g10:economic-models:other-things-equal}{other things equal}}.
\end{recall}

\section{From one buyer to the market}\label{sec:g10:demand:market}

Demand starts with each buyer. Ana and Ben take the bus in a town in Brenland, where the money is the dalen.
At a fare of $p$ dalens, Ana takes $q_A = 10 - 2p$ rides a week, as long as this is positive, and none if $p \ge 5$;
Ben takes $q_B = 12 - p$ rides, and none if $p \ge 12$. Ben lives further from his work and has fewer other ways to
get there.

\begin{definition}[Market demand]\label{def:g10:demand:market-demand}
The \emph{market demand}\index{market demand} for a good gives, at each price, the total quantity that all the buyers
in the market want to buy: the sum of the individual quantities demanded at that price.
\end{definition}

The \omterm{def:g10:demand:market-demand}{market demand} is found by adding quantities at the same price, not prices at the same quantity: on the graph,
the individual curves are added \emph{horizontally}. At a fare of 2 dalens, Ana takes 6 rides and Ben 10, so the market
takes 16. At a fare of 7, Ana takes none and Ben 5. Written out:
\[
  Q^d(p) = \begin{cases} (10 - 2p) + (12 - p) = 22 - 3p & \text{if } 0 \le p \le 5, \\
  12 - p & \text{if } 5 < p \le 12, \\ 0 & \text{if } p > 12. \end{cases}
\]

\begin{omfigure}
\begin{tikzpicture}[econ, x=0.38cm, y=0.42cm]
\draw[econ axes] (0,13.5) node[left, font=\small] {fare $p$} |- (24,0) node[below left=12pt and 0pt, font=\small] {rides a week};
\draw[demand, omDef!50] (0,5) -- (10,0);
\draw[demand, omDef!50] (0,12) -- (12,0);
\draw[curve] (0,12) -- (7,5) -- (22,0);
\draw[guide] (0,5) -- (7,5) -- (7,0);
\node[eqpoint] at (7,5) {};
\foreach \y in {5,12} \draw (0,\y) -- (-0.3,\y) node[left, font=\footnotesize] {\y};
\foreach \x in {7,10,12,22} \draw (\x,0) -- (\x,-0.3) node[below, font=\footnotesize] {\x};
\node[font=\footnotesize, text=omDef!70!black, anchor=south west] at (2.6,3.6) {Ana};
\node[font=\footnotesize, text=omDef!70!black, anchor=west] at (2.4,9.4) {Ben};
\node[font=\footnotesize, anchor=west] at (13,3.5) {market, $22 - 3p$};
\node[font=\footnotesize, anchor=south west] at (8.6,5.6) {kink: Ana stops riding};
\end{tikzpicture}

{\small Ana's and Ben's demands (light) and the \omterm{def:g10:demand:market-demand}{market demand} (black), added horizontally. Above a fare of 5 dalens
only Ben rides, and the market line is his; below it both ride, and the market line is flatter. The kink is at
$(7, 5)$.}
\end{omfigure}

\begin{proposition}[Adding straight demand lines]\label{prop:g10:demand:sum}
If every buyer's demand is a straight line down to zero, the \omterm{def:g10:demand:market-demand}{market demand} is made of straight pieces. On each piece
the slope, in quantity per dalen, is the sum of the slopes of the buyers who still buy at those prices. A kink appears
at each price at which a buyer stops buying.
\end{proposition}
\begin{proof}
At a price $p$, the market quantity is the sum of $a_i - b_i\,p$ over the buyers with $p < a_i/b_i$ (those who still
buy); the others contribute zero. Between two consecutive stopping prices the set of buyers does not change, so the sum
is $\sum a_i - \big(\sum b_i\big)\,p$, a straight line whose slope is the sum of their slopes. When $p$ crosses a stopping
price $a_j/b_j$, buyer $j$ drops out and the slope changes: a kink.
\end{proof}

With thousands of buyers, each stopping at a different price, the kinks are so close that the market curve looks smooth.
This is why a straight or gently curved \omterm{def:g10:demand:market-demand}{market demand} is a good description even though no single buyer is ``average''.

\section{The law of demand}\label{sec:g10:demand:law}

\begin{definition}[Law of demand]\label{def:g10:demand:law}
The \emph{law of demand}\index{law of demand} states that, \omterm{def:g10:economic-models:other-things-equal}{other things equal}, the quantity demanded of a good falls when
its price rises: the demand curve slopes downward.
\end{definition}

Two reasons work in the same direction. When gasoline becomes dearer \emph{relative to other things}, it pays to replace
some of it: a smaller car, a train, a bicycle, a house nearer work. And when a good that takes a large part of the budget
becomes dearer, the household is poorer, and buys less of every \omterm{def:g10:demand:normal}{normal good}, usually including that one. The law also holds for the
market even when some buyers do not respond at all: as long as some buyers stop buying, or buy less, when the price
rises, the market quantity falls.

\begin{model}{The law of demand}
\assumes \omterm{def:g10:economic-models:other-things-equal}{Other things equal}: the incomes of buyers, the prices of other goods, tastes, expectations and the number of
  buyers do not change while the price changes.
\predicts When the price of a good rises, the quantity bought falls; when it falls, the quantity bought rises.
\tested A randomised test: in an American health insurance experiment, families assigned at random to pay 95\,\% of
  their medical bills spent 845 dollars a year less on care than families whose care was free (grade 11). Consistent with it:
  after the real price of gasoline rose by about half between 1978 and 1981, oil use per person in the United States fell
  by almost a quarter by 1983, though a recession and new fuel-economy rules came at the same time.
\breaks Very poor households spending most of their budget on one staple food may buy more of it when its price rises
  (rice in Hunan, below).
% ledger: g11:hie-95
\end{model}
% figdata: oil_us (1978 3.22 -> 1981 4.75: +47.5 %; bpp 30.91 -> 1983 23.78: -23.1 %)

\begin{record}[Rice in Hunan: when a price fall cut purchases]
% ledger: g10:jensen-miller-2008
Among very poor households in Hunan, China, rice provided most of the calories. A study gave some of these households a
subsidy that lowered the price of rice. The \omterm{def:g10:demand:law}{law of demand} predicts that they would buy more rice. Many bought
\emph{less}: richer, thanks to the subsidy, they replaced part of their rice with meat and other foods they had been
unable to afford, while still getting enough calories. The demand curve of these households sloped upward over that
range, the behaviour economists call \emph{Giffen behaviour}. The same study found it more weakly for wheat in Gansu (Jensen and Miller 2008). It is the clearest case on record of an
exception to the \omterm{def:g10:demand:law}{law of demand}, and it needs exactly the conditions the model predicts: a staple that takes most of the
budget, and households too poor to buy enough of anything else.
\end{record}

The second chart puts the opening chart's numbers to a sterner test: year by year, does oil use fall when the price
rises? Out of 48 yearly changes, the two moved in opposite directions only 25 times. In 2009 and in 2020 both the
price and the quantity fell sharply. These were the years of a deep recession and of the first year of the pandemic:
incomes fell, travel stopped, and demand itself moved. The chart is not a demand curve; it is a cloud of points, each of
which mixes a move along the curve with a move of the curve.
% figdata: oil_summary (n 48, opposite 25); oil_changes (2009 -27.7, -4.6; 2020 -18.3, -11.8)

\begin{omchart}{Year-to-year changes in the real price of gasoline and in oil use per person, United States, 1977--2024,
in per cent. Each point is one year; the two labelled points are recession years in which both fell.}{US Bureau of Labor
Statistics; US Energy Information Administration; World Bank, World Development Indicators}
\begin{tikzpicture}
\begin{axis}[omscatter, width=0.78\linewidth, height=6.6cm, xmin=-32, xmax=38, ymin=-14, ymax=12,
  xlabel={change in the real price (\%)}, ylabel={change in oil use per person (\%)},
  extra x ticks={0}, extra y ticks={0}, extra tick style={grid style={black!60, thin}}]
\addplot[omDef, only marks, mark=*, mark size=1.5pt] table[x=dprice, y=dbpp] {figdata/grade-10/04-demand/out/oil_changes.csv};
\node[font=\footnotesize, anchor=south west] at (axis cs:-27.7,-4.6) {2009};
\node[font=\footnotesize, anchor=west] at (axis cs:-17.5,-11.8) {2020};
\end{axis}
\end{tikzpicture}
\end{omchart}

\section{What moves the curve}\label{sec:g10:demand:shifts}

The demand curve is drawn for given values of everything except the good's own price. When one of those other things
changes, the quantity demanded at \emph{every} price changes: the curve shifts. Five kinds of change matter.

\paragraph{Income.} For most goods, buyers buy more at every price when their income rises. Not for all: as incomes rise,
people take fewer long-distance buses and more flights, buy less of the cheapest bread and more of the dearer kinds.

\begin{definition}[Normal and inferior goods]\label{def:g10:demand:normal}
A good is \emph{normal}\index{normal good} when a rise in buyers' income, \omterm{def:g10:economic-models:other-things-equal}{other things equal}, raises the quantity demanded
at every price (the demand curve shifts right). It is \emph{inferior}\index{inferior good} when a rise in income lowers the
quantity demanded at every price.
\end{definition}
\phantomsection\label{def:g10:demand:inferior}

Whether a good is \omterm{def:g10:demand:normal}{normal} or \omterm{def:g10:demand:normal}{inferior} is a fact about buyers, not about the good: the same intercity bus can be \omterm{def:g10:demand:normal}{inferior}
for a well-paid office worker and \omterm{def:g10:demand:normal}{normal} for a student whose income has just risen.

\paragraph{Prices of related goods.} A rise in the price of train tickets raises the demand for bus rides; a rise in the
price of petrol lowers the demand for large cars.

\begin{definition}[Substitutes and complements]\label{def:g10:demand:substitutes}
Two goods are \emph{substitutes}\index{substitutes} when a rise in the price of one, \omterm{def:g10:economic-models:other-things-equal}{other things equal}, raises the demand
for the other: they can be used in each other's place. They are \emph{complements}\index{complements} when a rise in the
price of one lowers the demand for the other: they are used together.
\end{definition}
\phantomsection\label{def:g10:demand:complements}

\paragraph{Tastes and information.} A medical study linking a food to an illness shifts its demand left; a fashion shifts
the demand for what it favours right.

\paragraph{Expectations.} If buyers expect a price to rise next month, they buy more now: the demand curve of this month
shifts right. Queues at petrol stations before an announced tax rise are an example.

\paragraph{The number of buyers.} A town that grows, or a market that opens to foreign buyers, has more demand at every
price: the \omterm{def:g10:demand:market-demand}{market demand}, the horizontal sum of the individual ones, has more terms in its sum.

\begin{example}[A new rail line]\label{ex:g10:demand:rail}
A train line opens beside Ben's bus route, and Ben, at every fare, now takes 4 fewer bus rides a week: $q_B = 8 - p$. Ana's
demand does not change. The \omterm{def:g10:demand:market-demand}{market demand} becomes $18 - 3p$ for $p \le 5$ and $8 - p$ for $5 < p \le 8$. At a fare of 2 dalens,
rides fall from 16 to 12 without any change in the fare: the train is a substitute for the bus, and the demand for bus rides
has shifted left.
\end{example}

\section{Moving along the curve, and moving the curve}\label{sec:g10:demand:along}

The words matter. A change in the good's own price moves buyers \emph{along} the demand curve: the quantity demanded
changes, but demand --- the whole curve --- does not. A change in anything else moves the curve itself: demand changes.
``The price rose, so demand fell'' is a confusion; ``the price rose, so the quantity demanded fell'' is the \omterm{def:g10:demand:law}{law of demand}.

\begin{omfigure}
\begin{tikzpicture}[econ, x=0.55cm, y=0.5cm]
\begin{scope}
\draw[econ axes] (0,8.5) node[left, font=\small] {$p$} |- (9.5,0) node[below, font=\small] {$Q$};
\draw[demand] (0.5,7.5) -- (8.5,0.5);
\node[eqpoint] (a) at (2.5,5.75) {}; \node[eqpoint] (b) at (5.5,3.125) {};
\draw[phase arrow] (2.75,5.3) -- (5.2,3.2);
\node[font=\footnotesize, anchor=south west] at (2.6,5.8) {$p$ high};
\node[font=\footnotesize, anchor=south west] at (5.6,3.2) {$p$ low};
\node[font=\small\bfseries] at (4.7,-1.3) {a move along the curve};
\end{scope}
\begin{scope}[xshift=7.4cm]
\draw[econ axes] (0,8.5) node[left, font=\small] {$p$} |- (9.5,0) node[below, font=\small] {$Q$};
\draw[demand] (0.5,7.5) -- (6.5,0.5);
\draw[demand, shifted] (2.5,7.5) -- (8.5,0.5);
\draw[phase arrow] (2.3,4) -- (4.0,4);
\draw[guide] (0,4) -- (5.5,4);
\node[font=\footnotesize, anchor=west] at (6.0,6.5) {income rises};
\node[font=\small\bfseries] at (4.7,-1.3) {a shift of the curve};
\end{scope}
\end{tikzpicture}

{\small Left: a fall in the good's own price raises the quantity demanded along the same curve. Right: a rise in income
(for a \omterm{def:g10:demand:normal}{normal good}) shifts the whole curve to the right: at the same price, more is bought.}
\end{omfigure}

Now read the opening chart again. Between 1978 and 1983 the price rose and oil use fell, as the \omterm{def:g10:demand:law}{law of demand} predicts. But
other things did not stay equal: the recession of the early 1980s lowered incomes; cars were being built to new
fuel-economy rules; homes were heated with gas instead of oil. Part of the fall in oil use was a move along the curve, part
a shift of the curve. Between 1986 and 1998 the price was low, and oil use per person hardly rose: incomes grew, which
shifts demand right, while cars kept getting more efficient, which shifts it left. The record alone cannot separate the
two. How economists do it is a subject of the university volumes; the idea, already met in \cref{ch:g10:economic-models}, is
to find changes in the price that come from something that cannot have moved demand at the same time.

\begin{remark}[Two curves and one point]\label{rem:g10:demand:identification}
Each year's point on the opening chart is where that year's demand met that year's supply. If only supply moves, the points
trace out the demand curve. If only demand moves, they trace out the supply curve. When both move, the points are a cloud
that is neither. The next two chapters build the supply side and put the two together.
\end{remark}

\section{Exercises}

\begin{exercise}[$\star$]\label{exo:g10:demand:1}
A baker sells 120 loaves a day at 3 dalens, 100 at 4 dalens and 70 at 5 dalens. Is this consistent with the \omterm{def:g10:demand:law}{law of demand}?
By how many loaves does the quantity demanded fall for each dalen of price between 3 and 5?
\end{exercise}

\begin{exercise}[$\star$]\label{exo:g10:demand:2}
Say whether each pair is more likely to be \omterm{def:g10:demand:substitutes}{substitutes} or \omterm{def:g10:demand:substitutes}{complements} for most buyers: (a) tea and coffee; (b) printers and
ink cartridges; (c) bus rides and bicycles; (d) games consoles and games; (e) butter and margarine.
\end{exercise}

\begin{exercise}[$\star$]\label{exo:g10:demand:3}
A survey finds that households with higher incomes buy fewer kilograms of the cheapest rice and more of the dearer kinds,
at the same prices. What does this suggest about each kind of rice?
\end{exercise}

\begin{exercise}[$\star$]\label{exo:g10:demand:4}
Read the opening chart. In which year was the real price of gasoline highest? In which year was oil use per person lowest?
\end{exercise}

\begin{exercise}[$\star$]\label{exo:g10:demand:5}
Correct each sentence so that it uses ``demand'' and ``quantity demanded'' properly: (a) ``The price of cinema tickets fell,
so demand for cinema tickets rose.'' (b) ``Incomes rose, so the quantity demanded of restaurant meals rose, at the same
prices.''
\end{exercise}

\begin{exercise}[$\star\star$]\label{exo:g10:demand:6}
Two buyers have demands $q_1 = 20 - 4p$ (for $p \le 5$) and $q_2 = 15 - p$ (for $p \le 15$). Write the \omterm{def:g10:demand:market-demand}{market demand}, with its
pieces, and find the kink.
\end{exercise}

\begin{exercise}[$\star\star$]\label{exo:g10:demand:7}
Three identical buyers each have the demand $q = 6 - 2p$ (for $p \le 3$). Write the \omterm{def:g10:demand:market-demand}{market demand}. Compare its slope with one
buyer's, and explain the result with \cref{prop:g10:demand:sum}.
\end{exercise}

\begin{exercise}[$\star\star$]\label{exo:g10:demand:8}
For the demand for bus rides in a town, predict the direction of the shift (left, right, or none: a move along the curve)
for each event: (a) the bus fare rises; (b) the price of petrol rises; (c) a new car-sharing service starts; (d) the town's
population grows; (e) the bus company announces that fares will double next month (the demand for this month's monthly
passes).
\end{exercise}

\begin{exercise}[$\star\star$]\label{exo:g10:demand:9}
The \omterm{def:g10:demand:market-demand}{market demand} for rides is $Q = 22 - 3p$ for $p \le 5$ (Ana and Ben). A rise in Ana's income adds 2 rides a week to her
demand at every fare at which she rides. Write the new \omterm{def:g10:demand:market-demand}{market demand} for low fares, and the number of rides at a fare of 3
dalens before and after.
\end{exercise}

\begin{exercise}[$\star\star$]\label{exo:g10:demand:10}
In a town, the price of a cinema ticket was 8 dalens in 2024 and 9 dalens in 2025, and the number of tickets sold rose from
\num{50000} to \num{54000}. A new shopping centre with a large car park opened next to the cinema in 2025. Does this
contradict the \omterm{def:g10:demand:law}{law of demand}? Explain with the words ``along'' and ``shift''.
\end{exercise}

\begin{exercise}[$\star\star$]\label{exo:g10:demand:11}
% figdata: oil_changes (1980: 21.5, -8.7; 2021: 33.7, 9.2)
Read the second chart. In 1980 the real price of gasoline rose by 21.5\,\% and oil use per person fell by 8.7\,\%. In 2021
the price rose by 33.7\,\% and oil use rose by 9.2\,\%. Which year looks like a move along a demand curve? What, in 2021,
moved the demand curve, and in which direction?
\end{exercise}

\begin{exercise}[$\star\star\star$]\label{exo:g10:demand:12}
Describe an observation in a market that would contradict the \omterm{def:g10:demand:law}{law of demand}. Say what you would need to check before
concluding that the law had failed.
\end{exercise}

\begin{exercise}[$\star\star\star$]\label{exo:g10:demand:13}
Explain why the line joining the points of the opening chart, from 1978 to 1983, would not be the demand curve for oil
products in the United States. Would the true demand curve be steeper or flatter than that line, if incomes had fallen at
the same time? Explain.
\end{exercise}

\begin{exercise}[$\star\star\star$]\label{exo:g10:demand:14}
Explain, in words, why the households of the record box bought less rice when its price fell. What two conditions made this
possible, and why would one expect the same households to obey the \omterm{def:g10:demand:law}{law of demand} for meat?
\end{exercise}

\begin{exercise}[$\star\star\star$]\label{exo:g10:demand:15}
A market has two kinds of buyers of bread: \num{1000} poor households, for whom cheap bread is an \omterm{def:g10:demand:normal}{inferior good}, and \num{3000}
richer households, for whom it is \omterm{def:g10:demand:normal}{normal}. When all incomes rise by 10\,\%, each poor household buys 2 fewer loaves a month at
the old price and each richer household 1 more. What happens to the \omterm{def:g10:demand:market-demand}{market demand} at the old price? Is bread, for the market as
a whole, \omterm{def:g10:demand:normal}{normal} or \omterm{def:g10:demand:normal}{inferior}? What would change the answer?
\end{exercise}

\section{Problem: The Town's Bus Passes}

\begin{problem}[{Weekend problem --- three groups of riders, a fare rise and a new tram line}]
\label{pb:g10:demand:1}
In a town of Brenland, the bus company sells monthly passes. Its riders fall into three groups. At a price of $p$ dalens a
pass, the numbers of passes demanded a month are (each formula holds while it is positive, and the group buys none
otherwise):
\[
  \text{students: } q_S = 400 - 10p, \qquad \text{workers: } q_W = 900 - 15p, \qquad \text{retirees: } q_R = 300 - 10p .
\]
The price of a pass is 20 dalens.

\medskip
\noindent\textbf{Part I --- The groups.}
\begin{enumerate}
\item At what price does each group stop buying passes?
\item How many passes does each group buy at 20 dalens? What is the total?
\item Which group's demand falls by the most passes for each extra dalen?
\item The town's retirees receive a rise in their pensions, and their demand becomes $q_R = 360 - 10p$. Is a bus pass a \omterm{def:g10:demand:normal}{normal}
  or an \omterm{def:g10:demand:normal}{inferior good} for them? Before answering the next questions, cancel this change: the pensions do not rise.
\end{enumerate}

\medskip
\noindent\textbf{Part II --- The market.}
\begin{enumerate}[resume]
\item Write the \omterm{def:g10:demand:market-demand}{market demand} for $0 \le p \le 30$.
\item Write it for $30 < p \le 40$ and for $40 < p \le 60$.
\item Where are the kinks of the \omterm{def:g10:demand:market-demand}{market demand}? Explain each one in words.
\item Draw the \omterm{def:g10:demand:market-demand}{market demand}.
\end{enumerate}

\medskip
\noindent\textbf{Part III --- A fare rise.}
\begin{enumerate}[resume]
\item The company raises the price of a pass to 32 dalens. How many passes are sold now?
\item By how many passes, and by what percentage, does the number sold fall?
\item How many passes does each group buy at 32 dalens? Which group has left the market?
\item Compare the company's monthly takings at 20 and at 32 dalens.
\item Is the fall in passes sold a move along the demand curve or a shift of it?
\end{enumerate}

\medskip
\noindent\textbf{Part IV --- A tram line.}
\begin{enumerate}[resume]
\item A tram line opens past the university, and the students' demand for bus passes halves at every price: $q_S = 200 - 5p$.
  Are trams and buses \omterm{def:g10:demand:substitutes}{substitutes} or \omterm{def:g10:demand:substitutes}{complements} for the students?
\item Write the new \omterm{def:g10:demand:market-demand}{market demand} for $0 \le p \le 30$ and for $30 < p \le 40$.
\item With the tram line, how many passes are sold at 20 dalens, and how many at 32?
\item By how many passes, and by what percentage, does the fare rise reduce the number sold with the tram line?
\item Why does the same fare rise lose fewer passes with the tram line than without it?
\item By how many passes does the tram line alone reduce sales at 20 dalens? Is this a move along the curve or a shift?
\item State the result of the problem: the fall in passes sold after the fare rise from 20 to 32 dalens, without and with the
  tram line.
\end{enumerate}
\end{problem}
